\chapter{Quick Start}

This chapter gives a first complete GHATS workflow. The example follows the
RXTE/PCA walkthrough in the legacy manual: prepare an input list, produce a
PDS file, inspect it, extract a light curve, average the PDS, rebin it, and
plot it.

\section{Input Data}

For RXTE/PCA, the starting point is normally science-array or event data. The
supported data modes include binned data, single-bit data, event data, and
preprocessed GoodXenon data. The time and channel resolution depend on the
specific data mode.

The input to the PDS/FFT production program can be:

\begin{itemize}
\item a single science file;
\item a metafile, identified by a leading \texttt{@}, containing one input
  filename per line;
\item a meta-metafile, identified by a leading \texttt{@@}, containing one
  metafile name per line.
\end{itemize}

Files listed in the same metafile must be homogeneous: same mode, same time
resolution, same energy ranges, and same binning. For RXTE/PCA data, time
sorting is not needed by the production routine.

\begin{ghatsconsole}
$ ls FS3f* > @sb.lis
$ ls FS3b* > @event.lis
$ ls @sb.lis @event.lis > @@all.lis
\end{ghatsconsole}

\section{Producing a PDS}

Before running the production routine, compressed RXTE files and the required
auxiliary files should be uncompressed.

\begin{ghatsconsole}
$ gzip -d FS4a* FH5a*
\end{ghatsconsole}

Start GHATS and run the production routine:

\begin{ghatsconsole}
Ghats> gh_xte
\end{ghatsconsole}

When called without arguments, \routine{gh_xte} runs interactively. A file
selection window appears; in this example the selected file is
\texttt{@event.lis}. The program then shows the available channel ranges and
asks whether to use all of them or a selection.

\begin{ghatsnote}
If you are running IDL, \routine{gh_xte} can use the \keyword{TURBO} keyword,
which switches to a faster FFT engine from IDL\_EXTRACT by C. Markwardt.
\end{ghatsnote}

The routine asks whether to produce a PDS, stored as a power-spectrum file, or
an FFT, stored as a file of Fourier transforms. Unless the analysis requires
time lags or coherence, a PDS is usually the first product to make. In this
example the selected output type is \texttt{POWER}, and the output filename is
\texttt{event.pds}.

The routine can also rebin the data in time. In the example from the original
manual, the event data have a time resolution of 16 microseconds. Rebinning by
a factor of 4 reduces the highest frequency of interest while retaining enough
frequency range for high-frequency features.

The next choice is the number of time bins per transform, or equivalently the
duration of each uninterrupted data stretch. If a gap is reached before the
requested length is accumulated, the partial stretch is discarded and a new one
starts after the gap. This prevents gaps from entering individual Fourier
transforms. The routine accepts either the number of bins or the duration in
seconds; using a real number such as \texttt{32.0} requests a 32-second
transform.

\begin{ghatsconsole}
----------------------------------------------------------------
Program gh_xte     Tue Jan 3 15:26:00 2012     User: belloni
----------------------------------------------------------------
Working directory               :
/Users/belloni/data/xte/u1636/realtime/96087-01-84-10/pca
Input filename                  :
/Users/belloni/data/xte/u1636/realtime/96087-01-84-10/pca/@event.lis
Instrument                      : XTE/PCA
Observation ID                  : 96087-01-84-10
Source name                     : 4U_1636-53
Start date of observation       : 2011-12-29T04:01:36

Channels for accumulation       :    0-249
Type of output                  :    POWER
Output filename                 :    event.pds
Time resolution (s)             :    1/8192
Nyquist frequency (Hz)          :    4096.0000
Number of points per FFT        :    262144
Corresponding to length (s)     :    32.000000
Total number of FFTs            :    64
Time step                       :    32.000000
----------------------------------------------------------------
\end{ghatsconsole}

\section{Inspecting the Product}

The file \texttt{event.pds} contains a sequence of PDS, one per selected
uninterrupted data stretch. It also stores the corresponding count rate for
each transform and auxiliary information used by later GHATS routines.

The quickest overview is:

\begin{ghatsconsole}
Ghats> ghats_all,'event.pds'
GH_info ------------------------------------------------

PDS file name        :        event.pds

Observatory          :        XTE
Instrument           :        PCA
Target               :        4U_1636-53

N. of trafos    :   64
Trafo duration :    32.000000
N. of powers    :   262144
Start-end chans.:   0-249
Starting mjd    :   55924.168
Time step       :   0.031250000
Barycentered    :   NO
Spectral bins   :   3
Background      :   NO
------------------------------------------------------------
\end{ghatsconsole}

\routine{ghats_all} prints the basic header information and creates a summary
plot with the light curve and the averaged PDS. The same text information can
be obtained with:

\begin{ghatsconsole}
Ghats> gh_info,'event.pds'
\end{ghatsconsole}

\begin{ghatsnote}
Spaces between parts of IDL/GDL commands are harmless but not necessary. The
examples in this manual sometimes include spaces for readability.
\end{ghatsnote}

\section{Extracting and Plotting the Light Curve}

The light curve stored in the PDS file can be extracted and plotted with:

\begin{ghatsconsole}
Ghats> gh_licu,'event.pds',time,licu
Ghats> gh_plot_licu,time,licu
\end{ghatsconsole}

Here \param{time} and \param{licu} are arbitrary IDL/GDL variable names chosen
by the user.

\section{Averaging, Rebinning, and Plotting the PDS}

The workhorse routine for reading and averaging PDS files is \routine{ghx}.
With the default selections, the average PDS can be extracted as:

\begin{ghatsconsole}
Ghats> ghx,'event.pds',nu,pow,pow_e,/poisson
  64 power spectra selected
\end{ghatsconsole}

This reads the average PDS into three arrays: \param{nu} for frequency,
\param{pow} for power, and \param{pow_e} for the power error. The
\keyword{POISSON} keyword requests computation and subtraction of the
Poissonian component.

The PDS can then be rebinned:

\begin{ghatsconsole}
Ghats> gh_reb,nu,pow,pow_e,-100,x,y,ye
\end{ghatsconsole}

This writes the rebinned frequency, power, and error arrays to
\param{x}, \param{y}, and \param{ye}. Finally:

\begin{ghatsconsole}
Ghats> gh_plot_power,x,y,ye
\end{ghatsconsole}

\section{Custom PDS Selections}

The light curve can be used to select only transforms with count rates above a
chosen value. In the legacy example, the mean count rate is:

\begin{ghatsconsole}
Ghats> print,mean(licu)
  344.632
\end{ghatsconsole}

To average only PDS with count rates above this value:

\begin{ghatsconsole}
Ghats> ghx,'event.pds',nu,pow,pow_e,/poisson,rate=[344.632,10000]
  33 power spectra selected
\end{ghatsconsole}

Selections can also be combined with a time interval:

\begin{ghatsconsole}
Ghats> ghx,'event.pds',nu,pow,pow_e,/poisson,rate=[344.632,10000],time=[0,500]
  10 power spectra selected
\end{ghatsconsole}

or with an index interval:

\begin{ghatsconsole}
Ghats> ghx,'event.pds',nu,pow,pow_e,/poisson,index=[13,34]
  22 power spectra selected
\end{ghatsconsole}

\begin{ghatsnote}
GHATS command selections use indices starting from 1, not from 0.
\end{ghatsnote}

One can also pass an explicit selection array, for instance to select only odd
indices:

\begin{ghatsconsole}
Ghats> odd_indices = 2*indgen(32)+1
Ghats> ghx,'event.pds',nu,pow,pow_e,/poisson,rate=[344.632,10000], $
       time=[0,500],sel=odd_indices
  7 power spectra selected
\end{ghatsconsole}

\begin{ghatsnote}
The legacy Word/PDF extraction of this line appears malformed. The version
above gives the 1-based odd indices from 1 to 63.
\end{ghatsnote}

Finally, \routine{ghx} can return an rms-normalized PDS if the source and
background rates are supplied through the appropriate keyword. The background
rate must be estimated by the user.

\begin{ghatsconsole}
Ghats> ghx,'event.pds',nu,pow,pow_e,/poisson,rate=[344.632,10000], $
       time=[0,500],sel=odd_indices,rms=17.36
  7 power spectra selected
\end{ghatsconsole}

\begin{ghatsnote}
The scientific definitions and normalizations in this section are inherited
from the GHATS routines. This chapter is only a LaTeX conversion prototype; it
does not redefine the PDS, FFT, or normalization conventions.
\end{ghatsnote}
